\section{Tuning Secure Learned Bloom Filters}

This project was completed for CSC 592: Machine Learning Security.

\subsection{Background}

% AMQ filters, motivation for them + motivation for security
Database systems often require fast and small set data structures. For example, RocksDB~\cite{dong2021} reduces IO operations by filtering queries before 
reading from disk. As false positives do not impact correctness, they are permitted with bounded probability. Such approximate membership queries can be much
faster and more memory efficient than exact membership queries. However, adversaries may degrade performance by generating false positives. \\

\noindent
A Bloom filter~\cite{bloom1970} is a probabilistic data structure for approximate membership queries. Membership information is stored in an array of bits. To
query the filter, a key is hashed to several entries in the array. If all entries are set to 1 (i.e. the bitwise-AND of the entries is 1), then the query 
\emph{might} be in the set. False positives occur with probability $\epsilon$. Otherwise, the query \emph{definitely} is not in the set.  

% Bloom filter diagram
\begin{figure}[H]
    \centering
    \begin{tikzpicture}[>=Stealth, arraynode/.style={draw, minimum width=7mm, minimum height=6mm, outer sep=0, inner sep=0}]

        \node[arraynode, fill=lightyellow] (a1) {0};
        \node[arraynode, fill=lightorange, right=0 of a1] (a2) {1};
        \node[arraynode, fill=lightorange, right=0 of a2] (a3) {1};
        \node[arraynode, fill=lightyellow, right=0 of a3] (a4) {0};
        \node[arraynode, fill=lightyellow, right=0 of a4] (a5) {0};
        \node[arraynode, right=0 of a5, path picture={
        \shade[left color=lightorange, right color=lightred]
                (path picture bounding box.south west)
                rectangle (path picture bounding box.north east);
        }] (a6) {1};
        \node[arraynode, right=0 of a6, path picture={
        \shade[left color=lightred, right color=lightpink]
                (path picture bounding box.south west)
                rectangle (path picture bounding box.north east);
        }] (a7) {1};
        \node[arraynode, fill=lightyellow, right=0 of a7] (a8) {0};
        \node[arraynode, fill=lightyellow, right=0 of a8, path picture={
        \shade[left color=lightred, right color=lightpink]
                (path picture bounding box.south west)
                rectangle (path picture bounding box.north east);
        }] (a9) {1};
        \node[arraynode, fill=lightyellow, right=0 of a9] (a10) {0};
        \node[arraynode, fill=lightyellow, right=0 of a10, path picture={
        \shade[left color=lightpink, right color=lightpurple]
                (path picture bounding box.south west)
                rectangle (path picture bounding box.north east);
        }] (a11) {1};
        \node[arraynode, fill=lightyellow, right=0 of a11] (a12) {0};
        \node[arraynode, fill=lightpurple, right=0 of a12] (a13) {1};
        \node[arraynode, fill=lightyellow, right=0 of a13] (a14) {0};
        \node[arraynode, fill=lightyellow, right=0 of a14] (a15) {0};
        \node[arraynode, fill=lightpurple, right=0 of a15] (a16) {1};

        \node[circle, draw, above=12mm of a4, fill=lightorange] (x) {$x$};
        \draw[->] (x) to (a2.north);
        \draw[->] (x) to (a3.north);
        \draw[->] (x) to (a6.north);
        
        \node[circle, draw, above=12mm of a7, fill=lightred] (y) {$y$};
        \draw[->] (y) to (a6.north);
        \draw[->] (y) to (a7.north);
        \draw[->] (y) to (a9.north);

        \node[circle, draw, above=12mm of a10, fill=lightpink] (w) {$w$};
        \draw[->] (w) to (a7.north);
        \draw[->] (w) to (a9.north);
        \draw[->] (w) to (a11.north);
        
        \node[circle, draw, above=12mm of a13, fill=lightpurple] (z) {$z$};
        \draw[->] (z) to (a11.north);
        \draw[->] (z) to (a13.north);
        \draw[->] (z) to (a16.north);
        
    \end{tikzpicture}
    \caption{Bloom filter with 16-entry array and 3 hash functions. The filter contains keys $\mathcolorbox{lightorange}{x}$, $\mathcolorbox{lightred}{y}$, 
    $\mathcolorbox{lightpink}{w}$ and $\mathcolorbox{lightpurple}{z}$. Corresponding entries are marked by arrows and matching colors. For example, since the 
    filter contains $\mathcolorbox{lightorange}{x}$, the bitwise-AND 
    $\mathcolorbox{lightorange}{1} \mathbin{\&} \mathcolorbox{lightorange}{1} \mathbin{\&} \protect\gradientmathbox{lightorange}{lightred}{1}$ is $\mathcolorbox{lightorange}{1}$.}
    \label{fig:csc592_mlsec_bloom}
\end{figure}


% naor-eylon filters
\noindent
Secure Bloom Filters\footnote{Secure Bloom Filters are often (opaquely) called Naor-Eylon filters.} (SBFs)~\cite{} extend Bloom filters to provide strong
security guarantees. They apply a pseudorandom permutation\footnote{Pseudorandom permutations are not known to exist but are strongly suspected to.} to keys 
before querying the filter. If $\mathbf{P} \neq \mathbf{NP}$, there is no probabilistic polynomial time adversary that non-negligibly degrades the false 
positive rate. 

% learned bloom filter operation


% learned bloom filter diagram
\begin{figure}[H]
    \centering
\begin{tikzpicture}[>=Stealth, arraynode/.style={draw, minimum width=7mm, minimum height=6mm, outer sep=0, inner sep=0}]

\node[circle, draw, fill=lightorange] (x) {$x$};

\node[diamond, draw, right=8mm of x, aspect=2, fill=lightyellow, align=center] (Mx)
    {Learned Model\\$M(x) > \tau$?};

\node[rectangle, draw, below=8mm of Mx, fill=green!15, align=center] (Mxtrue)
    {$x \in S$};

\node[rectangle, draw, right=8mm of Mx, fill=lightyellow, align=center] (fn)
    {False Negative\\Backup Bloom Filter};

\node[rectangle, draw, right=8mm of fn, fill=red!20, align=center] (fnfalse)
    {$x \notin S$};

\node[rectangle, draw, fill=green!15, align=center] (fntrue) at (fn |- Mxtrue)
    {$x \in S$};

% Edges
\draw[->] (x) -- (Mx);
\draw[->] (Mx) -- node[right, font=\scriptsize] {Yes} (Mxtrue);
\draw[->] (Mx) -- node[above, font=\scriptsize] {No}  (fn);
\draw[->] (fn) -- node[above, font=\scriptsize] {No}  (fnfalse);
\draw[->] (fn) -- node[right, font=\scriptsize] {Yes} (fntrue);

\end{tikzpicture}
\caption{Learned Bloom filter. Learned model $M$ is queried with key $x$. If $M$ predicts membership with score greater than $\tau$, output that $x$ is
in the set. Otherwise, output the result of querying a Bloom filter containing the false negatives of $M$. False positives are possible but not false 
negatives.}
\label{fig:csc592_mlsec_lbf}
\end{figure}

% downtown bodgeda filter

% downtown bodgeda filter diagram
\begin{figure}[H]
    \centering
\begin{tikzpicture}[>=Stealth, arraynode/.style={draw, minimum width=7mm, minimum height=6mm, outer sep=0, inner sep=0}]

\node[circle, draw, fill=lightorange] (x) {$x$};
\node[diamond, draw, right=8mm of x, aspect=2, fill=lightyellow, align=center] (Mx)
    {Learned Model\\$M(x) > \tau$?};

\node[rectangle, draw, below=8mm of Mx, fill=lightyellow, align=center] (tp)
    {True Positive\\Backup Bloom Filter};
\node[rectangle, draw, right=8mm of tp, fill=red!20, align=center] (tpfalse)
    {$x \notin S$};
\node[rectangle, draw, below=8mm of tp, fill=green!15, align=center] (tptrue)
    {$x \in S$};

\node[rectangle, draw, right=8mm of Mx, fill=lightyellow, align=center] (fn)
    {False Negative\\Backup Bloom Filter};
\node[rectangle, draw, right=8mm of fn, fill=red!20, align=center] (fnfalse)
    {$x \notin S$};
\node[rectangle, draw, below=8mm of fn, fill=green!15, align=center] (fntrue)
    {$x \in S$};

% Edges
\draw[->] (x)  -- (Mx);
\draw[->] (Mx) -- node[right, font=\scriptsize] {Yes} (tp);
\draw[->] (Mx) -- node[above, font=\scriptsize] {No}  (fn);
\draw[->] (tp) -- node[above, font=\scriptsize] {No}  (tpfalse);
\draw[->] (tp) -- node[right, font=\scriptsize] {Yes} (tptrue);
\draw[->] (fn) -- node[above, font=\scriptsize] {No}  (fnfalse);
\draw[->] (fn) -- node[right, font=\scriptsize] {Yes} (fntrue);

\end{tikzpicture}
\caption{Learned Bloom filter. Learned model $M$ is queried with key $x$. If $M$ predicts membership with score greater than $\tau$, output the result of 
querying a Bloom filter containing true positives of $M$. Otherwise, output the result of querying a Bloom filter containing the false negatives of $M$. False 
positives are possible but not false negatives.}
\label{fig:csc592_mlsec_lbf}
\end{figure}

% research questions
% optimal backup filter sizing algorithm
% empirical attack on DBF

Achieving optimal performance from LBFs requires tuning several parameters including a decision threshold $\tau$, and backup filter false positive rate(s) to 
attain the desired false positive rate $\epsilon$. Parameter tuning also impacts the filter's memory footprint and, for DBFs, worse-case positive rate. Tuning 
procedures for LBFs do not directly translate to DBFs, as they do not balance the tradeoff between expected and worst-case false positive rates. Therefore, I 
ask:
\begin{enumerate}
    \item How can the decision threshold and backup filter false positive rates be chosen to minimize SLBF memory footprint given desired expected and maximum
    false positive rates?
    \item 
\end{enumerate}

\subsection{Summary}

% RQ1:

% description of tuning algorithm (math)

% other part (dp)

% RQ2:

% dbf implementation

% dataset

% empirical attack + results

TODO: A 400–500-word summary of the technical contributions of the work, what was
accomplished overall, and the limitations of their work.

\subsection{Reflection}

TODO: A 250-word self reflection of what skills were developed or enhanced in completing those
projects.

% attending talks

% understanding theory papers

% rapid ideation and development